%% kapa: exam=1 round=34 year=2023 period=1 form=A subject=civil_law subject_ko=민법 seq=31 number=31
%% source: Q-Net 감정평가사 기출문제 (cst00309) artlSeq=5211789 file=raw/1008/5211789/2023년도 제34회 감정평가사 1차시험 1교시.pdf
%% conversion: kordoc 4.24.51 -> tools/gen1.py (auto); review=proofread-sonnet; flags=none
\begin{question}{31}
공동소유에 관한 설명으로 옳지 않은 것은? (다툼이 있으면 판례에 따름)
\begin{choices}
  \item 공유자는 다른 공유자의 동의없이 공유물을 처분하거나 변경하지 못한다.
  \item 합유는 수인이 조합체로서 물건을 소유하는 형태이고, 조합원은 자신의 지분을 조합원 전원의 동의없이 처분할 수 없다.
  \item 합유물에 대한 보존행위는 합유자 전원의 동의를 요하지 않는다.
  \item 구조상·이용상 독립성이 있는 건물부분이라 하더라도 구분소유적 공유관계는 성립할 수 없다.
  \item 공유물분할 금지약정은 갱신할 수 있다.
\end{choices}
\end{question}
\answer{4}
